% ===================================================================
%  TABELLA nr. 13 — L'hotel. --- Il restaurant. Il chaffè.
%  Traduction 2026 d'apres texto/io, controlee sur texto/fr, texto/ca
%  et texto/oc. Consigne : voir texto/rm/10-tabella-01.tex.
%
%  LE TABLEAU 12 A DONNE AU CRITERE SA PREMIERE REUSSITE ; CELUI-CI
%  LUI DONNE SON PREMIER ECHEC, ET L'ECHEC APPREND PLUS.
%
%  Le critere des tableaux 11 et 12 disait : qui fonde nomme. Le
%  chemin de fer, monte par une compagnie du pays, s'appelle
%  « Viafier Retica » ; la caisse d'epargne, fondee au village,
%  s'appelle « banca da spargn » ; les pompiers, corps venu du
%  dehors, s'appellent « pumpiers ».
%
%  Or l'hotellerie de montagne est l'invention des Grisons. C'est en
%  Engadine, dans une vallee romanche, qu'on a ouvert la saison
%  d'hiver ; les hoteliers etaient du pays, l'argent aussi. Par le
%  critere, ce vocabulaire-la devrait etre indigene.
%
%  IL NE L'EST PAS, ET PAS D'UN SEUL MOT : « hotel », « portier »,
%  « camerier », « buffet », « menu », « reception », « grum »,
%  « paquebot » — et le pourboire de l'alinea 9 est « buonamaun »,
%  le beau-main italien.
%
%  POURQUOI LA DIFFERENCE ? Elle ne tient ni au fondateur ni a la
%  date, elle tient a QUI DOIT COMPRENDRE LE MOT. Les voyageurs de la
%  Viafier Retica etaient des gens du pays qui allaient au marche ;
%  les clients de l'hotel etaient anglais et allemands. UN
%  VOCABULAIRE SUIT SA CLIENTELE, PAS SON FONDATEUR.
%
%  Le critere du tableau 11 se corrige donc ici, comme le tableau 15
%  luxembourgeois avait corrige le tableau 15 lituanien : il ne
%  suffit pas de demander qui a bati la maison, il faut demander a
%  qui elle ouvre sa porte.
%
%  ET LE POURBOIRE FAIT LA TROISIEME REPONSE DE LA SERIE A LA MEME
%  PIECE DE MONNAIE : le tcheque et le lituanien CALQUENT le
%  Trinkgeld allemand ; le luxembourgeois le PREND tel quel ; le
%  romanche prend le « buonamano » italien. Trois langues, trois
%  sources, et chaque fois la source est celle d'ou vient LE
%  SERVICE, non celle d'ou vient la frontiere.
% ===================================================================

%%K t13-noto-1 noto
\VUnotes{143.2mm}{%
\textit{\VUgras{(*) Per ils detagls da la chombra, vesair la tabella nr.~6.}}%
}

%%K t13-apar-1 apar
\VUsaut{3.0mm}
\VUtitre{15.0pt}{60}{TABELLA nr. 13}
\VUsaut{1.6mm}
\VUpk{10.2pt}{40}{L'hotel. --- Il restaurant.}
\VUsaut{2.0mm}
\VUpk{10.2pt}{40}{Il chaffè.}
\VUsaut{3.4mm}

%%K t13-01-1 p
1. --- Essend arrivà ier saira a Royan, sun jau ì a lozar a l'«\textit{Hotel da
l'Ocean}», ch'in ami m'aveva recumandà.

%%K t13-01-2 p
Ins m'ha dà ina bella \VUgras{chombra} \textsuperscript{(2)}, sin il terz \VUgras{plaun}
\textsuperscript{(1)}, nua ch'jau hai durmì fitg bain (*).

%%K t13-02-1 p
2. --- Questa damaun, sortind da mia chombra, en la quala la \VUgras{chombrera}
\textsuperscript{(3)} m'aveva purtà l'ensolver, sun jau entrà en la \VUgras{chombra da
fimar} \textsuperscript{(5)}, che vegn era duvrada sco \VUgras{chombra da lectura}
\textsuperscript{(4)} per leger las \VUgras{gasettas} \textsuperscript{(6)}, fimond ina
\VUgras{cigaretta} \textsuperscript{(8)}. Suenter avair legì sun jau ì giu cun
l'\VUgras{ascensur} \textsuperscript{(9)} e sun ì a spassegiar tras la citad.

%%K t13-03-1 p
3. --- Perquai ch'jau aveva emblidà da dumandar l'\VUgras{impiegà} \textsuperscript{(12)}
da l'hotel a tge ura ch'ins dat ils pasts a la \VUgras{maisa cuminaivla}
\textsuperscript{(11)}, sun jau turnà baud ed hai profità per ir a ma bognar en la
\VUgras{stanza da bogn} \textsuperscript{(10)} da l'hotel. Lura sun jau ì danovamain al
\VUgras{biro da la cassa} \textsuperscript{(15)} per telefonar ad in da mes amis. Ma il
\VUgras{telefon} \textsuperscript{(13)} era occupà; vesend mia embarazza, ha la
\VUgras{cassiera} \textsuperscript{(14)}, ch'era per registrar in \VUgras{quint}
\textsuperscript{(17)} sin il \VUgras{register dals giasts} \textsuperscript{(16)}, rugà il
\VUgras{portier} \textsuperscript{(18)} da trametter il \VUgras{grum} \textsuperscript{(19)} a far
mia commissiun \VUgras{cun il velo} \textsuperscript{(20)}.

%%K t13-04-1 p
4. --- Jau l'hai ingrazià e sun sortì per dar ina dunativa al \VUgras{purtader}
\textsuperscript{(24)}, che m'aveva purtà da la staziun sin ses \VUgras{char da maun}
\textsuperscript{(25)} mia \VUgras{chascha da chapels} \textsuperscript{(26)}, \VUgras{mia
valisch} \textsuperscript{(27)} e \VUgras{mes satg da viadi} \textsuperscript{(28)}. Il
\VUgras{purtader da brevs} \textsuperscript{(21)}, ch'era gist arrivà, m'ha dà ina
\VUgras{brev} \textsuperscript{(22)}.

%%K t13-05-1 p
5. --- Jau sun restà anc in pau sin il sulier da la porta, sper la \VUgras{chascha da
brevs} \textsuperscript{(23)}, guardond la circulaziun sin la via. Il \VUgras{conductur}
\textsuperscript{(30)} da l'\VUgras{omnibus} \textsuperscript{(29)} chargiava sin ses char la
\VUgras{crappa} \textsuperscript{(31)} d'in \VUgras{viagiader} \textsuperscript{(32)}, ch'il
\VUgras{patrun da l'hotel} \textsuperscript{(33)} accumpagnava. In giuven \VUgras{telegrafist}
\textsuperscript{(35)} purtava senza pressa in \VUgras{telegram} \textsuperscript{(34)} per in
client da l'hotel; in \VUgras{turist} \textsuperscript{(36)} cun ses \VUgras{apparat da
fotografar} \textsuperscript{(38)} en bandulera consultava ses \VUgras{cudesch da guid}
\textsuperscript{(37)}.

%%K t13-tit-1 sub
\VUsaut{4.0mm}
\VUpk{10.2pt}{40}{L'ura dal gentar.}
\VUsaut{3.0mm}

%%K t13-06-1 p
6. --- Avant il grigl durmigliava in \VUgras{commissiunari} \textsuperscript{(39)} senza
quità tranter sia \VUgras{hocca da purtar} \textsuperscript{(40)} e sia \VUgras{chascha da
lustrar} \textsuperscript{(41)}, entant ch'ina giuvna
\VUgras{lavandera} \textsuperscript{(42)},
che purtava in \VUgras{chanaster} \textsuperscript{(43)} da biancaria, rieva pervi da la
mina solemna dal \VUgras{maistral da maisa} \textsuperscript{(44)}, ch'era en pes sper la
porta.

%%K t13-07-1 p
7. --- Vesend il \VUgras{cuschinunz} \textsuperscript{(45)}, ch'è sortì da sia
\VUgras{cuschina} \textsuperscript{(47)} situada en il \VUgras{sutterren} \textsuperscript{(48)}
per trametter in \VUgras{giuven da servetsch} \textsuperscript{(46)} a far ina commissiun,
sun jau ma algordà ch'l'ura dal gentar è prest, e jau sun entrà en il restaurant,
traversond la curt en la quala in \VUgras{chauffeur} \textsuperscript{(50)} garava ses
\VUgras{auto} \textsuperscript{(49)}, entant ch'in \VUgras{mecanicist} \textsuperscript{(52)}
reparava, tenor las indicaziuns dal \VUgras{ciclist} \textsuperscript{(53)}, in
\VUgras{triciclo a petroli} \textsuperscript{(51)} ch'aveva gist gì in accident.

%%K t13-08-1 p
8. --- Jau hai dumandà in giuven \VUgras{grum} \textsuperscript{(54)}, che purtava
\VUgras{magiels da biera} \textsuperscript{(55)}, sche la maisa cuminaivla saja sut la
veranda e, pervi da sia resposta affirmativa, hai jau prendì plaz vi d'ina da las
maisas ed hai laschà dar a mai in past excellent.

%%K t13-noto-2 noto
\VUnotes{143.2mm}{%
\textit{\VUgras{(*) Cun agid da la glista dals plats introducida en il vocabulari
vegn il magister a pudair variar facilmain il menu che suonda.}}%
}

%%K t13-tit-2 sub
\VUsaut{4.0mm}
\VUpk{10.2pt}{40}{In excellent gentar.}
\VUsaut{3.0mm}

%%K t13-09-1 p
9. --- Il camerier m'ha purtà il menu (*) dal gentar e la glista dals vins. Jau hai
tschernì ed ordinà sco \VUgras{hors-d'oeuvre crevettas} cun \VUgras{paintg} e, sco
\VUgras{entrada}, \VUgras{tripa a la moda da Caen}. Jau n'hai betg maglià
\VUgras{pesch}, ma hai dumandà ina porziun da \VUgras{cossa} d'agnè e, sco \VUgras{plat
da verduras}, \VUgras{spinat} cun roma. Lura, perquai che nagin dals \VUgras{intermezs}
na ma plascheva, hai jau laschà purtar differents desserts, \VUgras{chaschiel},
\VUgras{fritgas} e \VUgras{biscuits}. Jau hai plinavant agiuntà in excellent
\VUgras{vin} da Saint-Emilion e, fitg cuntent da mes past, hai jau bandunà la maisa
suenter avair dà al \VUgras{camerier} \textsuperscript{(57)} in bel \VUgras{buonamaun}
\textsuperscript{(59)}.

%%K t13-10-1 p
10. --- Vesend ch'jau era pront a partir, m'ha l'\VUgras{administratur}
\textsuperscript{(60)} proponì d'ir sin il balcun per admirar il splendid panorama da
l'\VUgras{Ocean} \textsuperscript{(62)}. Dalunsch vaseva ins \VUgras{barcuttas}
\textsuperscript{(63)} da pestga, \VUgras{navs a vela} \textsuperscript{(64)} ed in enorm
\VUgras{paquebot} \textsuperscript{(65)}, dal qual la nera siluetta sortiva sin ils alvs
\VUgras{nivels} \textsuperscript{(67)} da l'\VUgras{orizont} \textsuperscript{(66)}. Ina brisa
lediera faseva flottar la \VUgras{bandiera} \textsuperscript{(70)} fitgada sur
l'\VUgras{ensaina}
\textsuperscript{(69)} da la \VUgras{halla} \textsuperscript{(68)}.

%%K t13-tit-3 sub
\VUsaut{3.0mm}
\VUpk{10.2pt}{40}{Divertiments.}
\VUsaut{3.0mm}

%%K t13-11-1 p
11. --- Il spectacul era uschè interessant ch'jau hai decidì da muntar damaun sin la
terrassa dal chaffè cun ina \VUgras{binocla} \textsuperscript{(71)} u in \VUgras{cannocchial}
\textsuperscript{(72)}.

%%K t13-12-1 p
12. --- Bandunond il balcun, sun jau ì en il chaffè da l'hotel per prender ina
\VUgras{bavranda} \textsuperscript{(73)}, giugond ina \VUgras{partida da bigliard}
\textsuperscript{(75)}. Senza ma fermar hai jau traversà la sala dal chaffè, en la quala
blers consuments giugavan \VUgras{schah} \textsuperscript{(78)}, \VUgras{dama}
\textsuperscript{(79)}, \VUgras{dominos} \textsuperscript{(80)}, \VUgras{triktrak}
\textsuperscript{(81)} u \VUgras{cartas} \textsuperscript{(82)}, e sun ì directamain si en la
\VUgras{sala da bigliard} \textsuperscript{(74)}.

%%K t13-13-1 p
13. --- Cur ch'la partida è stada finida, sun jau ì giu al plaunterren ed hai dumandà
dal camerier in \VUgras{cuvert} \textsuperscript{(84)},
\VUgras{palpiri} \textsuperscript{(83)},
ina \VUgras{marca postala} \textsuperscript{(85)}, ina \VUgras{penna} \textsuperscript{(87)}
ed
\VUgras{encaust} \textsuperscript{(86)} per scriver ina brev.

%%K t13-13-2 p
Lura hai jau clamà in \VUgras{vitturin} \textsuperscript{(92)}, dal qual il
\VUgras{chalesch} \textsuperscript{(88)} era fermà avant il chaffè. Jau al hai dumandà il
pretsch tenor las uras u per ina cursa e, cur ch'nus essan stads d'accord davart il
pretsch, m'ha el avert la \VUgras{porta dal chalesch} \textsuperscript{(89)} e m'ha
installà. El è puspè muntà sin il \VUgras{sez} \textsuperscript{(91)}, ha fatg partir
svelt ils chavals cun ina \VUgras{scurriada} \textsuperscript{(93)}, e nus essan partids per
ina charmanta spassegiada.

\VUsaut{6.0mm}
\VUfilet{22mm}
